% Native reconstruction of source slides 3--45.
\section[Matrix operations]{Matrix Operations}

% Slide 3
\sectioncard{2.1}{Matrix Operations}

% Slide 4
\begin{frame}{Matrix Terminology}
\begin{block}{Definition}
A \textbf{matrix} is an ordered rectangular array of elements.
\end{block}
\[
A=\begin{bmatrix}
a_{11}&a_{12}&\cdots&a_{1n}\\
a_{21}&a_{22}&\cdots&a_{2n}\\
\vdots&\vdots&\ddots&\vdots\\
a_{m1}&a_{m2}&\cdots&a_{mn}
\end{bmatrix}.
\]
\begin{itemize}
  \item $A$ has $m$ rows and $n$ columns, so its size is $m\times n$.
  \item The scalar in row $i$, column $j$ is $a_{ij}$, the $(i,j)$-entry.
  \item We write $A=A_{m\times n}=(a_{ij})_{m\times n}=(a_{ij})$.
\end{itemize}
\end{frame}

% Slide 5
\begin{frame}{Matrix Sizes: Examples}
\centering
\renewcommand{\arraystretch}{1.25}
\begin{tabular}{c@{\qquad}c}
$\begin{bmatrix}1&0&3&5\\-9&6&4&3\end{bmatrix}$ & $2\times4$\\[3ex]
$\begin{bmatrix}2&3&5&9\end{bmatrix}$ & $1\times4$\\[3ex]
$\begin{bmatrix}1\\2\\4\end{bmatrix}$ & $3\times1$\\[3ex]
$\begin{bmatrix}4\end{bmatrix}$ & $1\times1$
\end{tabular}
\end{frame}

% Slide 6
\begin{frame}{Two Ways to Describe a Matrix}
For an $m\times n$ matrix $A$, we may work with its entries:
\[
A=\begin{bmatrix}
a_{11}&\cdots&a_{1j}&\cdots&a_{1n}\\
\vdots&&\vdots&&\vdots\\
a_{i1}&\cdots&a_{ij}&\cdots&a_{in}\\
\vdots&&\vdots&&\vdots\\
a_{m1}&\cdots&a_{mj}&\cdots&a_{mn}
\end{bmatrix},
\]
or with its columns:
\[
A=[\,\vect a_1\ \vect a_2\ \cdots\ \vect a_n\,].
\]
\end{frame}

% Slide 7
\begin{frame}{Some Special Matrices}
\begin{itemize}
  \item A \textbf{square matrix of order $n$} has size $n\times n$.
  \[
  \begin{bmatrix}13&6&2\\2&2&2\\2&2&2\end{bmatrix}
  \]
  \item A \textbf{row matrix} (row vector) of order $n$ has the form
  \[A=(a_1,a_2,\ldots,a_n).\]
  \item A \textbf{column matrix} (column vector) of order $n$ has the form
  \[B=\begin{bmatrix}a_1\\a_2\\\cdots\\a_n\end{bmatrix}.\]
\end{itemize}
\end{frame}

% Slide 12
\begin{frame}{Matrix Operations}
  Today, we are going to formally introduce the matrix operations:
\begin{itemize}
  \item \textcolor{lecturemagenta}{"+": Addition}
  \item \textcolor{lecturemagenta}{"c": Scalar multiplication}
  \item \textcolor{lecturemagenta}{"*": Matrix multiplication}
  \item \textcolor{lecturemagenta}{"T": Transpose of a matrix}
  \item \textcolor{lecturemagenta}{"$^{-1}$" or "/": Inverse of a matrix}
\end{itemize}
\end{frame}

% Slide 13
\begin{frame}{Sums and Scalar Multiples}

What does linear system preserve?\pause
\begin{itemize}
  \item \textcolor{lecturemagenta}{Addition}
  \item \textcolor{lecturemagenta}{Scalar multiplication}
\end{itemize}\pause
\vfill
Vectors satisfy this rule. If we write vectors into a matrix, matrix satisfies the rule too.


\end{frame}

% Slide 14
\begin{frame}{Recall: Vector Arithmetic}
Let $\vect u,\vect v\in\R^n$.
\begin{itemize}
  \item $\vect u=\vect v$ if and only if their corresponding entries are equal.
  \item $\vect u+\vect v$ is obtained by adding corresponding entries.
  \item $c\vect u$ is obtained by multiplying every entry of $\vect u$ by the scalar $c$.
\end{itemize}
\end{frame}

% Slide 15
\begin{frame}{Same-Size and Equal Matrices}
Now, for matrices:

Two matrices have the \textbf{same size} when they have the same number of rows and the same number of columns. For example,
\[
\begin{bmatrix}1&2\\5&6\\3&7\end{bmatrix}
\quad\text{and}\quad
\begin{bmatrix}14&3\\8&4\\3&9\end{bmatrix}
\]
both have size $3\times2$.
\medskip

Two matrices are \textbf{equal} when they have \textbf{1)} the same size and \textbf{2)} all corresponding entries are equal.
\end{frame}

% Slide 16

% Slide 10
\begin{frame}{Matrix Addition: Zero Matrices}
What we need to define first, if we want to define an addition?\pause

The "0"!\pause

An $m\times n$ matrix whose entries are all zero is a \textbf{zero matrix}, written $0_{m\times n}$ or simply $0$.
\[
\begin{bmatrix}
0&0&0&0\\
0&0&0&0\\
0&0&0&0\\
0&0&0&0
\end{bmatrix}
\ne
\begin{bmatrix}0&0&0&0\end{bmatrix}.
\]

It adds any matrix equals the matrix itself. $0+A=A$.\pause

The size of a zero matrix is usually clear from context. You know the size if you know the size of A, don't you?
\end{frame}


\begin{frame}{Matrix Addition}
If $A=(a_{ij})_{m\times n}$ and $B=(b_{ij})_{m\times n}$ have the same size, then
\[
A+B=(a_{ij}+b_{ij})_{m\times n}
=\begin{bmatrix}
a_{11}+b_{11}&\cdots&a_{1n}+b_{1n}\\
a_{21}+b_{21}&\cdots&a_{2n}+b_{2n}\\
\vdots&\ddots&\vdots\\
a_{m1}+b_{m1}&\cdots&a_{mn}+b_{mn}
\end{bmatrix}.
\]
\alert{The sum $A+B$ is defined only when $A$ and $B$ have the same size.}
\end{frame}

% Slide 17
\begin{frame}{Matrix Addition: Example}
\small
\begin{align*}
&\begin{bmatrix}12&3&-5\\1&-9&0\\3&6&8\end{bmatrix}
+\begin{bmatrix}1&8&9\\6&5&4\\3&2&1\end{bmatrix}\\[1ex]
&=\begin{bmatrix}12+1&3+8&-5+9\\1+6&-9+5&0+4\\3+3&6+2&8+1\end{bmatrix}
=\begin{bmatrix}13&11&4\\7&-4&4\\6&8&9\end{bmatrix}.
\end{align*}
\end{frame}

% Slide 18
\begin{frame}{Scalar Multiplication}
If $\lambda$ is a scalar and $A=(a_{ij})$, then
\[
\lambda A=(\lambda a_{ij})
=\begin{bmatrix}
\lambda a_{11}&\lambda a_{12}&\cdots&\lambda a_{1n}\\
\lambda a_{21}&\lambda a_{22}&\cdots&\lambda a_{2n}\\
\vdots&\vdots&\ddots&\vdots\\
\lambda a_{m1}&\lambda a_{m2}&\cdots&\lambda a_{mn}
\end{bmatrix}.
\]
Equivalently, if $A=[\,\vect a_1\ \cdots\ \vect a_n\,]$, then
\[
\lambda A=[\,\lambda\vect a_1\ \cdots\ \lambda\vect a_n\,].
\]
\end{frame}

% Slide 19
\begin{frame}{Negation and Subtraction}
We have learned $O$. What operation makes $O$?\pause

$a+0=0$

$a-a=0$ \pause

For a matrix, we do not need to invent 'minus".

When performing \textbf{scalar multiplications}, negation means multiply by -1.

We define $-A=(-1)A$, so every entry multiplies by -1. We then write
\[
A-B=A+(-1)B.
\]
\begin{exampleblock}{Questions}
\begin{enumerate}
  \item For what matrix $B$ does $A+B=A$?
  \item What is $A-A$? \pause $O$
  \item Is $A-A=B-B$ always true for matrices $A$ and $B$?
\end{enumerate}
\end{exampleblock}
\end{frame}


% Slide 21
\begin{frame}{Sums and Scalar Multiples: Example}
Let
\[
A=\begin{bmatrix}2&4&-5\\0&7&-2\end{bmatrix},\qquad
B=\begin{bmatrix}1&-3&3\\0&2&-4\end{bmatrix}.
\]
Then
\[
A+B=\begin{bmatrix}3&1&-2\\0&9&-6\end{bmatrix},
\qquad
2(A+B)=\begin{bmatrix}6&2&-4\\0&18&-12\end{bmatrix}.
\]
\end{frame}



% Slide 8
\begin{frame}{Scalar multiplication and Diagonal Matrix}
A \textbf{diagonal matrix} is a square matrix whose nondiagonal entries are zero:
\[
A=\operatorname{diag}(\lambda_1,\lambda_2,\ldots,\lambda_n)
=\begin{bmatrix}
\lambda_1&0&\cdots&0\\
0&\lambda_2&\cdots&0\\
\vdots&\vdots&\ddots&\vdots\\
0&0&\cdots&\lambda_n
\end{bmatrix}.
\]
\pause
The diagonal entries need not all be zero, and some may be zero. We love this form of matrix. \pause

If diagonal entries are 1, it is an identity matrix $I$.
\end{frame}


\begin{frame}{Scalar multiplication and Diagonal Matrix}
We love this type of matrix because we can multiply a matrix by row or by column.\pause
\[
A=\operatorname{diag}(\lambda_1,\lambda_2,\ldots,\lambda_n)
=\begin{bmatrix}
\lambda_1&0&\cdots&0\\
0&\lambda_2&\cdots&0\\
\vdots&\vdots&\ddots&\vdots\\
0&0&\cdots&\lambda_n
\end{bmatrix}.
\]
\pause

\begin{itemize}
  \item If A is diagonal,\pause
  \item Multiplication at left: $AB$ scales $\lambda_i$ for each ith row.\pause
  \item Multiplication at right: $BA$ scales $\lambda_i$ for each ith column.
\end{itemize}

A scalar scales everything equally.

A diagonal matrix lets us scale different directions differently.
\end{frame}
\begin{frame}{Essence of Matrix}
  A table of numbers $\rightarrow$ operation or transformation.
\end{frame}

\begin{frame}{Scalar Multiplication and Diagonal Matrix}
  Diagonal Matrix bridges \textbf{scalar} multiplication with \textbf{matrix} multiplication!\pause

  For a $2 \times 2$ matrix $A$:

  \[
cA
=\begin{bmatrix}
c&0\\
0&c
\end{bmatrix}A=(cI)A.
\]
  
\end{frame}

% Slide 20
\begin{frame}{Properties of Sums and Scalar Multiples}
\theorembox{Theorem 1}{Let $A,B,C$ be matrices of the same size, and let $r,s$ be scalars. Then}
\begin{columns}[T]
\column{.48\textwidth}
\begin{align*}
A+B&=B+A,\\
(A+B)+C&=A+(B+C),\\
A+0&=A.
\end{align*}
\column{.48\textwidth}
\begin{align*}
r(A+B)&=rA+rB,\\
(r+s)A&=rA+sA,\\
r(sA)&=(rs)A.
\end{align*}
\end{columns}
\end{frame}


% Slide 9
\begin{frame}{Matrix Multiplication}

What do we need to define first, if we want to define an multiplication?\pause

The "1", or identity matrix!\pause

The \textbf{identity matrix} of order $n$ is the diagonal matrix
\[
I=I_n=\operatorname{diag}(1,1,\ldots,1)
=\begin{bmatrix}
1&0&\cdots&0\\
0&1&\cdots&0\\
\vdots&\vdots&\ddots&\vdots\\
0&0&\cdots&1
\end{bmatrix}.
\]
All main diagonal entries are $1$.
\end{frame}


% Slide 22
\begin{frame}{Matrix--Vector Multiplication}
If $A=[\,\vect a_1\ \vect a_2\ \cdots\ \vect a_n\,]$ is $m\times n$ and $\vect x\in\R^n$, then
\[
A\vect x=
[\,\vect a_1\ \vect a_2\ \cdots\ \vect a_n\,]
\begin{bmatrix}x_1\\x_2\\\vdots\\x_n\end{bmatrix}
=x_1\vect a_1+x_2\vect a_2+\cdots+x_n\vect a_n.
\]
Thus $A\vect x$ is a linear combination of the columns of $A$, with entries of $\vect x$ as weights.
\end{frame}

% Slide 23
\begin{frame}[shrink=4]{Matrix--Vector Multiplication: Example}
Compute $A\vect x$, where
\[
A=\begin{bmatrix}2&3&4\\-1&5&-3\\6&-2&8\end{bmatrix},\qquad
\vect x=\begin{bmatrix}x_1\\x_2\\x_3\end{bmatrix}.
\]
Then
\begin{align*}
A\vect x
&=x_1\begin{bmatrix}2\\-1\\6\end{bmatrix}
+x_2\begin{bmatrix}3\\5\\-2\end{bmatrix}
+x_3\begin{bmatrix}4\\-3\\8\end{bmatrix}\\
&=\begin{bmatrix}
2x_1+3x_2+4x_3\\
-x_1+5x_2-3x_3\\
6x_1-2x_2+8x_3
\end{bmatrix}.
\end{align*}
\end{frame}

% Slide 24
\begin{frame}{Row-by-Row View of $A\vect x$}
Each entry of $A\vect x$ is a row-vector product:
\[
\begin{array}{c}
\begin{bmatrix}2&3&4\end{bmatrix}\vect x=2x_1+3x_2+4x_3,\\[2ex]
\begin{bmatrix}-1&5&-3\end{bmatrix}\vect x=-x_1+5x_2-3x_3,\\[2ex]
\begin{bmatrix}6&-2&8\end{bmatrix}\vect x=6x_1-2x_2+8x_3.
\end{array}
\]
\end{frame}

% Slide 25
\begin{frame}{Matrix Multiplication}
Let
\[
A=(a_{ij})_{m\times r},\qquad B=(b_{ij})_{r\times n}.
\]
The product $AB$ is the $m\times n$ matrix whose $(i,j)$-entry is
\[
(AB)_{ij}=a_{i1}b_{1j}+a_{i2}b_{2j}+\cdots+a_{ir}b_{rj}
=\sum_{k=1}^{r}a_{ik}b_{kj}.
\]
\end{frame}

% Slide 26
\begin{frame}{The Row--Column Rule}
\centering
\begin{tikzpicture}[>=Stealth]
\node (a) {$A_{m\times r}$};
\node[right=3.2cm of a] (b) {$B_{r\times n}$};
\node[right=3.2cm of b] (c) {$(AB)_{m\times n}$};
\draw[->,thick] (a)--(b);
\draw[->,thick] (b)--(c);
\node[below=1.1cm of b,draw,rounded corners,inner sep=8pt] {$\displaystyle (AB)_{ij}=\underbrace{\text{row }i\text{ of }A}_{1\times r}\,\underbrace{\text{column }j\text{ of }B}_{r\times1}$};
\end{tikzpicture}
\vfill
\[
(AB)_{ij}=a_{i1}b_{1j}+a_{i2}b_{2j}+\cdots+a_{ir}b_{rj}.
\]
\end{frame}

% Slide 27
\begin{frame}[shrink=8]{Matrix Multiplication: Examples}
\[
\begin{bmatrix}-2&4\\1&-2\end{bmatrix}
\begin{bmatrix}2&4\\-3&-6\end{bmatrix}
=\begin{bmatrix}-16&-32\\8&16\end{bmatrix}.
\]
\medskip
If
\[
A=\begin{bmatrix}1&0&-1&2\\-1&1&3&0\\0&5&-1&4\end{bmatrix},\quad
B=\begin{bmatrix}0&3&4\\1&2&1\\3&1&-1\\-1&2&1\end{bmatrix},
\]
then
\[
AB=\begin{bmatrix}-5&6&7\\10&2&-6\\-2&17&10\end{bmatrix}.
\]
Here $A$ is $3\times4$, $B$ is $4\times3$, and $AB$ is $3\times3$.
\end{frame}

% Slide 28
\begin{frame}{When Is a Product Defined?}
The product $AB$ exists if and only if the number of columns of $A$ equals the number of rows of $B$.
\[
\underbrace{A_{m\times r}}_{\text{$r$ columns}}
\underbrace{B_{r\times n}}_{\text{$r$ rows}}
=C_{m\times n}.
\]
For example, a $3\times3$ matrix cannot multiply a $2\times3$ matrix on its right.
\medskip

A $1\times1$ matrix is usually written as a scalar:
\[
\begin{bmatrix}1&2&3\end{bmatrix}
\begin{bmatrix}3\\2\\1\end{bmatrix}
=1\cdot3+2\cdot2+3\cdot1=10.
\]
\end{frame}




% Slide 30
\begin{frame}{A Product May Exist in Only One Order}
Matrix multiplication shares some properties with real-number multiplication, but not all.
\medskip

Let
\[
A=\begin{bmatrix}1&5\\-2&0\end{bmatrix},\qquad
B=\begin{bmatrix}1&0&-6\\-2&1&-1\end{bmatrix}.
\]
Then
\[
AB=\begin{bmatrix}-9&5&-11\\-2&0&12\end{bmatrix},
\]
while $BA$ is not defined because the inner dimensions do not match.
\end{frame}

\begin{frame}{Multiplication May Not Be Swappable!!}
  \begin{itemize}
    \item We all know $a \times b = b \times a$.
    \item But this is special. \pause
    \item Think about \[\mathbf{v}=
    \begin{bmatrix}1\\2\end{bmatrix}
    \]
    \item Scale x twice, then swap x and y.
    \item Swap x and y, then scale x twice. Are they the same?
  \end{itemize}
  
\end{frame}


\begin{frame}{Multiplication May Not Be Swappable!!}
  \pause
  \[\mathbf{D}=\begin{bmatrix}2&0\\0&1\end{bmatrix}, \mathbf{S}=\begin{bmatrix}0&1\\1&0\end{bmatrix}\]\pause

  \[\mathbf{SD}=\begin{bmatrix}0&1\\2&0\end{bmatrix}, \mathbf{DS}=\begin{bmatrix}0&2\\1&0\end{bmatrix}\]\pause

  \[\mathbf{SD} \neq \mathbf{DS}\]
\end{frame}

% Slide 31
\begin{frame}{Matrix Multiplication Not Commutative}
For
\[
A=\begin{bmatrix}-1&2&3\\0&3&-2\end{bmatrix},\qquad
B=\begin{bmatrix}2&1\\0&2\\-1&0\end{bmatrix},
\]
both products exist, but
\[
AB=\begin{bmatrix}-5&3\\2&6\end{bmatrix},\qquad
BA=\begin{bmatrix}-2&7&4\\0&6&-4\\1&-2&-3\end{bmatrix}.
\]
Thus $AB\ne BA$ in general.
\end{frame}

% Slide 32
\begin{frame}{Zero Products}
Let
\[
A=\begin{bmatrix}1&1\\0&0\end{bmatrix},\qquad
B=\begin{bmatrix}1&0\\-1&0\end{bmatrix}.
\]
Then
\[
AB=\begin{bmatrix}0&0\\0&0\end{bmatrix},\qquad
BA=\begin{bmatrix}1&1\\-1&-1\end{bmatrix}.
\]
Therefore
\[
AB=0\centernot\implies A=0\text{ or }B=0,
\qquad
A\ne0,\ B\ne0\centernot\implies AB\ne0.
\]
\end{frame}


\begin{frame}{Matrix Multiplication is Applying Operations}

\begin{block}{Key idea}
A matrix acts on a vector.  
In a product of matrices, the matrix \textbf{closest to the vector acts first}.
\end{block}

\[
AB\mathbf{x} = A(B\mathbf{x})
\]

\[
\mathbf{x}
\;\xrightarrow{\;B\;}\;
B\mathbf{x}
\;\xrightarrow{\;A\;}\;
A(B\mathbf{x}) = AB\mathbf{x}
\]

\vspace{0.5em}

\[
BA\mathbf{x} = B(A\mathbf{x})
\]

\[
\mathbf{x}
\;\xrightarrow{\;A\;}\;
A\mathbf{x}
\;\xrightarrow{\;B\;}\;
B(A\mathbf{x}) = BA\mathbf{x}
\]
\end{frame}
\begin{frame}{Matrix Multiplication is Applying Operations}

\begin{alertblock}{Takeaway}
When applying transformation to a vectors to the right, the \textbf{rightmost} matrix acts first.  
So \(AB\mathbf{x}\) and \(BA\mathbf{x}\) usually mean applying operations in a different order, and therefore in general
\[
AB \ne BA.
\]
\end{alertblock}

\end{frame}
% Slide 29
\begin{frame}{Properties of Matrix Multiplication}
\theorembox{Theorem 2}{Let $A$ be $m\times n$ and let $B,C$ have sizes for which the indicated operations are defined. Then}
\begin{align*}
A(BC)&=(AB)C &&\text{associative law},\\
A(B+C)&=AB+AC &&\text{left-distributive law},\\
(B+C)A&=BA+CA &&\text{right-distributive law},\\
r(AB)&=(rA)B=A(rB),\\
I_mA&=A=AI_n.
\end{align*}
\end{frame}

% Slide 33
\begin{frame}{Exercise: Row and Column Products}
Let
\[
\vect a=\begin{bmatrix}1&-1&2\end{bmatrix},\qquad
\vect b=\begin{bmatrix}2\\3\\0\end{bmatrix}.
\]
Then
\[
\vect a\vect b=1\cdot2+(-1)\cdot3+2\cdot0=-1,
\]
whereas
\[
\vect b\vect a=
\begin{bmatrix}
2&-2&4\\3&-3&6\\0&0&0
\end{bmatrix}.
\]
Again, matrix multiplication is not commutative.
\end{frame}

% Slide 34
\begin{frame}{Powers of a Matrix}
For a square matrix $A$ and a natural number $k$,
\[
A^k=\underbrace{AA\cdots A}_{k\text{ factors}}.
\]
For natural numbers $k$ and $\ell$,
\[
(A^k)^\ell=A^{k\ell},\qquad
A^{k+\ell}=A^kA^\ell.
\]
In general,
\[
(AB)^k\ne A^kB^k,
\]
because the factors cannot be rearranged unless $AB=BA$.
\end{frame}

% Slide 35
\begin{frame}[shrink=4]{The Identity Matrix}
When $k=0$, $A^0$ is interpreted as the identity matrix. The $n\times n$ identity matrix is
\[
I=(\delta_{ij}),\qquad
\delta_{ij}=\begin{cases}1,&i=j,\\0,&i\ne j.\end{cases}
\]
For example,
\[
I_3=\begin{bmatrix}1&0&0\\0&1&0\\0&0&1\end{bmatrix}.
\]
If $A=\begin{bmatrix}3&4&1\\2&6&3\\0&1&8\end{bmatrix}$, then $I_3A=A=AI_3$.
\begin{block}{Property}
An identity matrix behaves like the number $1$ in multiplication.
\end{block}
\end{frame}

% Slide 36
\begin{frame}[shrink=6]{Exercises: True or False?}
\begin{enumerate}
  \item $A^2=A$ implies $A=I$ or $A=0$. \alert{False.}
  \[
  A=\begin{bmatrix}1&-1&1\\1&-1&1\\1&-1&1\end{bmatrix}
  \quad\text{satisfies}\quad A^2=A,
  \]
  yet $A\ne I$ and $A\ne0$.
  \item $AC=BC$ and $C\ne0$ imply $A=B$. \alert{False in general.}
  \item $A^2-B^2=(A+B)(A-B)$. \alert{False in general,} because
  $(A+B)(A-B)=A^2-AB+BA-B^2$.
\end{enumerate}
\end{frame}

\begin{frame}{Transpose}
\begin{columns}
  \begin{column}{.3\linewidth}
    \[
\begin{bmatrix}a&b&c\\d&e&f\end{bmatrix}
\xrightarrow{\text{transpose}}
\begin{bmatrix}a&d\\b&e\\c&f\end{bmatrix}
\]
  \end{column}\pause
  \begin{column}{.7\linewidth}
  \begin{figure}
    \includegraphics[width=.4\linewidth]{penguin.png}
    \includegraphics[width=.4\linewidth]{penguintranspose.png}
  \end{figure}  
  \end{column}
\end{columns}
\end{frame}
% Slide 37
\begin{frame}{Transpose: Main Diagonal Entries}
For $A=(a_{ij})_{m\times n}$, the diagonal entries are
\[
a_{11},a_{22},a_{33},\ldots,
\]
and they form the \textbf{main diagonal} of $A$.
\begin{center}
\[
\begin{bmatrix}
\color{red}a_{11}&a_{12}&\cdots&a_{1n}\\
a_{21}&\color{red}a_{22}&\cdots&a_{2n}\\
\vdots&\vdots&\color{red}\ddots&\vdots\\
a_{m1}&a_{m2}&\cdots&\color{red}a_{mn}
\end{bmatrix}
\]
\end{center}
\end{frame}

% Slide 38
\begin{frame}{Transpose of a Matrix}
The transpose of an $m\times n$ matrix $A$, denoted by $A^T$, is the $n\times m$ matrix whose columns are the corresponding rows of $A$.
\[
A=\begin{bmatrix}a_{11}&a_{12}&\cdots&a_{1n}\\a_{21}&a_{22}&\cdots&a_{2n}\\\vdots&\vdots&\ddots&\vdots\\a_{m1}&a_{m2}&\cdots&a_{mn}\end{bmatrix},
\quad
A^T=\begin{bmatrix}a_{11}&a_{21}&\cdots&a_{m1}\\a_{12}&a_{22}&\cdots&a_{m2}\\\vdots&\vdots&\ddots&\vdots\\a_{1n}&a_{2n}&\cdots&a_{mn}\end{bmatrix}.
\]
For example,
\[
\begin{bmatrix}1&3\\2&5\\7&4\end{bmatrix}^{T}
=\begin{bmatrix}1&2&7\\3&5&4\end{bmatrix}.
\]
\end{frame}

% Slide 39
\begin{frame}[shrink=5]{Properties of the Transpose}
If $A=(a_{ij})$ and $A^T=(b_{ij})$, then $b_{ij}=a_{ji}$. Equivalently, rows become columns in the same order.
\medskip

\theorembox{Theorem 3}{For matrices of appropriate sizes and any scalar $r$:}
\begin{align*}
(A^T)^T&=A,\\
(A+B)^T&=A^T+B^T,\\
(rA)^T&=rA^T,\\
(AB)^T&=B^TA^T.
\end{align*}
The order reverses when a product is transposed.
\end{frame}

% Slide 40
\begin{frame}{Computing a Transpose of a Product: Way 1}
Let
\[
A=\begin{bmatrix}2&0&-1\\1&3&2\end{bmatrix},\qquad
B=\begin{bmatrix}1&7&-1\\4&2&3\\2&0&1\end{bmatrix}.
\]
First multiply:
\[
AB=\begin{bmatrix}0&14&-3\\17&13&10\end{bmatrix}.
\]
Then transpose:
\[
(AB)^T=\begin{bmatrix}0&17\\14&13\\-3&10\end{bmatrix}.
\]
\end{frame}

% Slide 41
\begin{frame}{Computing a Transpose of a Product: Way 2}
Use $(AB)^T=B^TA^T$ directly:
\[
(AB)^T=
\begin{bmatrix}1&4&2\\7&2&0\\-1&3&1\end{bmatrix}
\begin{bmatrix}2&1\\0&3\\-1&2\end{bmatrix}
=\begin{bmatrix}0&17\\14&13\\-3&10\end{bmatrix}.
\]
The same identity gives
\[
(A^2)^T=(A^T)^2,\qquad (A^m)^T=(A^T)^m.
\]
\end{frame}

% Slide 42
\begin{frame}[shrink=9]{Exercise: Transpose and Products}
Let
\[
A=\begin{bmatrix}-1&2&1\\0&-1&2\end{bmatrix},\quad
B=\begin{bmatrix}1&0&3\\2&1&-1\end{bmatrix},\quad
C=\begin{bmatrix}-1&1&2\\3&-2&1\\0&0&4\end{bmatrix}.
\]
Compute $(A+2B)C$ and $(2A-B)C^T$.
\[
(A+2B)C=
\begin{bmatrix}1&2&7\\4&1&0\end{bmatrix}C
=\begin{bmatrix}5&-3&32\\-1&2&9\end{bmatrix},
\]
\[
(2A-B)C^T=
\begin{bmatrix}-3&4&-1\\-2&-3&5\end{bmatrix}C^T
=\begin{bmatrix}5&-18&-4\\9&5&20\end{bmatrix}.
\]
\end{frame}

% Slide 43
\begin{frame}{Symmetric and Skew-Symmetric Matrices}
\begin{block}{Symmetric matrix}
A matrix $A$ is symmetric if $A^T=A$. Equivalently, $A$ is square and $a_{ji}=a_{ij}$.
\end{block}
\[
\begin{bmatrix}2&6&1\\6&8&0\\1&0&6\end{bmatrix}
\quad\text{is symmetric.}
\]
\begin{block}{Skew-symmetric matrix}
A matrix $A$ is skew-symmetric if $A^T=-A$.
\end{block}
\end{frame}

% Slide 44
\begin{frame}{Examples of Symmetry}
\[
A=\begin{bmatrix}1&2&1\\2&0&-1\\1&-1&2\end{bmatrix}
\quad\text{is symmetric,}
\]
\[
B=\begin{bmatrix}0&2&1\\-2&0&-1\\-1&1&0\end{bmatrix}
\quad\text{is skew-symmetric.}
\]
Indeed, $A^T=A$ and $B^T=-B$.
\medskip

For every skew-symmetric matrix,
\[
a_{ii}=0,\qquad a_{ji}=-a_{ij}\quad(i\ne j).
\]
\end{frame}

% Slide 45
\begin{frame}[shrink=7]{Exercises: Symmetric Matrices}
\begin{enumerate}
  \item If $A$ and $B$ are symmetric, show that $A+B$ is symmetric.
  \item Show that $AA^T$ and $A^TA$ are symmetric:
  \[
  (AA^T)^T=(A^T)^TA^T=AA^T,
  \qquad
  (A^TA)^T=A^T(A^T)^T=A^TA.
  \]
  \item If $A,B$ are symmetric and $k$ is a scalar, determine which matrices are always symmetric:
  \[
  \text{(A) }kAB,\qquad
  \text{(B) }AB^T,\qquad
  \text{(C) }A+kB^T,\qquad
  \text{(D) }kA-B^T.
  \]
\end{enumerate}\pause
The answers are (C) and (D).
\end{frame}


% Slide 11
\begin{frame}{Special Matrices: Recap}
For $A=(a_{ij})_{m\times n}$:
\begin{description}
  \item[Square matrix] $m=n$.
  \item[Row or column matrix] One dimension equals $1$.
  \item[Identity matrix] $a_{ii}=1$ and $a_{ij}=0$ for $i\ne j$.
  \item[Diagonal matrix] $a_{ij}=0$ for $i\ne j$.
  \item[Zero matrix] Every entry is zero.
\end{description}
\end{frame}

